% Figure block for the R-Universe complete-chain manuscript.
% The figures are generated by code/make_figures_and_tables.py.

\begin{figure}[H]
\centering
\includegraphics[width=0.86\linewidth]{fig01_expansion_Ez.pdf}
\caption{Expansion history $E(z)=H(z)/H_0$ for the Pantheon+ full covariance + DESI DR2 BAO + cosmic-chronometer best fits. R1 remains close to \LCDM but shifts the preferred matter density and introduces a controlled exponent $\theta$.}
\label{fig:expansion}
\end{figure}

\begin{figure}[H]
\centering
\includegraphics[width=0.86\linewidth]{fig02_om_diagnostic.pdf}
\caption{Om diagnostic. In flat \LCDM, ${\rm Om}(z)$ is constant and equals $\Omega_{\rm m0}$. The R1 and R2 branches predict a redshift drift, providing a direct background-level null test.}
\label{fig:om}
\end{figure}

\begin{figure}[H]
\centering
\includegraphics[width=0.98\linewidth]{fig03_desi_dr2_bao_pulls.pdf}
\caption{DESI DR2 BAO diagnostic pulls. The plotted values use diagonal errors only for visual readability; the numerical likelihood uses the full DESI covariance matrix.}
\label{fig:bao}
\end{figure}

\begin{figure}[H]
\centering
\includegraphics[width=0.86\linewidth]{fig04_pantheon_residuals.pdf}
\caption{Binned Pantheon+ residuals after analytic marginalization over the additive supernova offset. The R1 improvement is distributed over redshift rather than caused by a single outlying point.}
\label{fig:snres}
\end{figure}

\begin{figure}[H]
\centering
\includegraphics[width=0.96\linewidth]{fig05_model_selection.pdf}
\caption{Information-criterion summary. Negative values favor R-Universe over \LCDM. R1 is AIC-favored in the DESI DR2 and DES-Dovekie configurations but is not BIC-favored.}
\label{fig:modelselection}
\end{figure}

\begin{figure}[H]
\centering
\includegraphics[width=0.82\linewidth]{fig06_growth_Dz.pdf}
\caption{Minimal GR growth solution normalized to $D(0)=1$. This is a closed late-time consistency calculation, not a substitute for a full Boltzmann implementation.}
\label{fig:growth}
\end{figure}
